\chapter{Matching both weak C2 packets to the four strong packets}
\label{component:p8b53b68d8af6}
\noindent\textbf{Component }\path{p8b53b68d8af6}\par
{\small\noindent Source: \path{CORE/prior/inherited/c2_matching_strength.tex}\par}
{\small\noindent SHA256: \path{8b53b68d8af6a308d041e96a3bf10b86a1c991010146e95da907ad2d40e15ec0}\par}
{\small\noindent Inventory use: Exact complete-parent reflection, both product-bound corner joins, and full ordinary-angle radial source proof\par}
\medskip
\section{Matching both weak C2 packets to the four strong packets}
\label{p8b53b68d8af6:sec:c2-matching-strength}

A coefficient condition symmetric in the two weak packets closes both
commuting angle corners at once.  It also gives a second exact radial
source family whose contact angles have ordinary positive residue.

\begin{theorem}[Two matched weak packets]
\label{p8b53b68d8af6:thm:c2-matching-strength}
Let the original physical C2 parent have $d=1$, $\lambda=\beta=0$,
and event weights $(n,m,C,A,D,\alpha)$.  Suppose $C,\alpha>0$ are
infinitesimal, all four other weights are positive, and each of
\[
 Cn,\ Cm,\ CA,\ CD,\qquad
 \alpha n,\ \alpha m,\ \alpha A,\ \alpha D
\]
has bounded reciprocal in the coefficient field.  Then the full parent
is tensor SOS at every independent projective test and every real or
complex state.
\end{theorem}
\begin{proof}
Lemma~\ref{pc7dacaa3258e:lem:c2-four-strong-outer} applies directly to every commuting
residue test except the finite zero-zero corner.
For that corner use the exact reflection
\begin{equation}
 (u,v;W,X,Y,Z;n,m,C,A,D,\alpha)
 \longmapsto
 (1/v,1/u;Z,X,Y,W;n,D,\alpha,A,m,C).
\label{p8b53b68d8af6:eq:c2ms-reflection}
\end{equation}
More explicitly, if $F$ denotes the full parent with its indicated
arguments, then
\begin{equation}
 F(u,v;W,X,Y,Z;n,m,C,A,D,\alpha)
 =u^2v^2F(1/v,1/u;Z,X,Y,W;n,D,\alpha,A,m,C).
\label{p8b53b68d8af6:eq:c2ms-reflection-identity}
\end{equation}
This is a homogeneous polynomial identity, so its interpretation in
reciprocal charts includes literal zero and infinity endpoints.

For a short verification, the reflected $R_2,R_5$ amplitudes, after
multiplication by $uv$, are $-R_5,-R_2$.  The reflected $C$ and
$\alpha$ packets exchange, while the $n,A$ packets are fixed.
If $\mathcal H$ is the signed core, its common expanded form is
\begin{align*}
 \mathcal H/2={}&u^2(W^2+Z^2)+(1+v^2+u^2v^2)(X^2+Y^2)\\
 &+2u^2v(WX-YZ)+2u(XZ-WY)-2uv(WZ+XY).
\end{align*}
Multiplying its reflected expression by $u^2v^2$ fixes the first,
second and last terms, and exchanges the middle two.  This proves the
whole-parent identity, including its signed part.

The reflected weak weight is $\alpha$, its four strong weights are
$n,D,A,m$, and its other nonnegative weak packet is $C$.
Thus all hypotheses of Lemma~\ref{pc7dacaa3258e:lem:c2-four-strong-outer} hold for
the reflected parent.  Its infinity-infinity chart is precisely the
original zero-zero chart.  Restore the homogeneous pure-testing
factors to obtain the certificate there as well.

Every noncommuting residue test is covered by the already-established
physical-source theorem, after positive pure-coefficient normalization
if a coefficient is unbounded.  Undoing that normalization preserves
SOS.  These alternatives exhaust the tests, and real/imaginary splitting
handles complex states.  No coefficient scale was compared with an
independent testing scale.
\end{proof}

\begin{corollary}[Radial C2 source with ordinary contact angles]
\label{p8b53b68d8af6:cor:c2-ordinary-angle-radial}
Use the same rational source \eqref{p9f1235511111:eq:c2u-source}, now with source
half-angle $c$ a positive bounded unit.  Suppose $h=b-a^2t,r,g,j$ are
positive infinitesimals and $a,1-r,t$ are positive bounded units.
Then the full parent is tensor SOS at every independent projective test
and every real or complex state.
\end{corollary}
\begin{proof}
Put $z=c^2$ and continue to write $C=e_3$ for the weak event weight.
The exact source identities
\[
 b-1=-r+g-zh,\qquad t-1=zr-j,\qquad
 a-1=\frac{g+j-(1+z)(r+h)}{t(a+1)}
\]
remain valid for this ordinary source angle.  Thus $a,b,t$ have residue
one.  The positive source identity
\[
 N=(1+z)(b+r)h+a^2\{g+j+rt(1+a^2z)\}
\]
shows that $N>0$ is infinitesimal and $h/N,r/N,g/N,j/N$ are bounded:
each of the four positive coefficients in this sum is a bounded unit.

Here is the entire normalized coefficient calculation.  Write
$z_0=\overline z>0$.  Substitution of $a=b=t=1,r=0$ in the bounded
source numerators and unit denominators gives
\begin{gather}
 \overline\Delta=2(1+z_0)^2,\qquad
 \overline{C/N}=\frac1{4(1+z_0)^2},\qquad
 \overline{\alpha/C}=z_0^2,\notag\\
 \overline{rn}=\overline{hA}=\frac1{4(1+z_0)},\qquad
 \overline{gm}=\overline{jD}=\frac14.
\label{p8b53b68d8af6:eq:c2ms-ordinary-residues}
\end{gather}
For example, the exact quotients used in the first line are
\[
 \frac CN=\frac{1+zr}{2a^2t\Delta},\qquad
 \frac\alpha C=\frac{z^2(b+r)a^2t}{(1-zh)(1+zr)}.
\]
The four strong normalized products use the numerators of
\eqref{p9f1235511111:eq:c2u-source} after multiplication by $r,h,g,j$, respectively;
no vanishing source gap is inverted in this calculation.
All six weights are therefore positive, $C,\alpha$ are infinitesimal,
and their normalized quotients displayed above are bounded units.
For instance
\[
 \frac1{Cn}=\frac{r/N}{(C/N)(rn)},\qquad
 \frac1{\alpha n}=\frac{C/\alpha}{Cn};
\]
the other six product bounds follow in the same way from $g/N,h/N,j/N$.
These are exactly the hypotheses of Theorem~\ref{p8b53b68d8af6:thm:c2-matching-strength}.

For completeness, the source is physical by the unchanged exact contact
argument.  The rational stationary identities
\eqref{p9f1235511111:eq:c2u-contact} are valid at this value of $c$ as well.
At that noncommuting contact the state matrix has a nonzero kernel
vector.  Its restriction to $Z=0$ is positive definite: the signed core
is nonnegative there and the positive event packets have zero common
kernel.  Hence the other three eigenvalues are positive.  The exact
contact and angular stationarity give the stationary zero of the
Laurent determinant; its constant-plus-positive-exponentials form makes
this its global minimum.  The determinant criterion and homogeneous
continuity prove physicality, exactly as for the small-angle source.
Theorem~\ref{p8b53b68d8af6:thm:c2-matching-strength} now applies.

This corollary concerns the stated radial subfamily of the ordinary
source-angle cell.  It does not assert that every ordinary-angle C2
source has the required matching strengths or enters this subfamily.
\end{proof}
